\documentclass[10pt,a4paper]{article}

\usepackage{geometry}
\geometry{
  left=15mm,
  right=15mm,
  top=15mm,
  bottom=15mm
}
\usepackage{amsmath}
\usepackage{amssymb}
\usepackage{enumitem}

\setlist[enumerate,1]{label=(\arabic*),leftmargin=2.2em}
\setlist[enumerate,2]{label=(\alph*),leftmargin=2.2em}

\begin{document}

\begin{center}
  {\Large \textbf{The University of Tokyo 2023}}\\[2mm]
  {\large \textbf{Problems 1--3}}
\end{center}

\section*{Problem 1}

Answer the following questions.

\begin{enumerate}
  \item The function $f(x,y)$ with real variables $x,y$ is defined as follows:
        \[
          f(x,y)=
          \begin{vmatrix}
            1 & x_1 & y_1\\
            1 & x_2 & y_2\\
            1 & x & y
          \end{vmatrix}.
        \]
        Show that the set of solutions of the equation $f(x,y)=0$ is a line passing through two points $(x_1,y_1)$, $(x_2,y_2)$ on the $xy$ plane, where $x_1\neq x_2$.

  \item Find the value of the determinant
        \[
          \begin{vmatrix}
            1 & x_1 & x_1^2\\
            1 & x_2 & x_2^2\\
            1 & x_3 & x_3^2
          \end{vmatrix}
        \]
        in factored form.

  \item Show that there is a unique curve
        \[
          y=a_0+a_1x+a_2x^2
        \]
        passing through three points $(x_1,y_1)$, $(x_2,y_2)$, $(x_3,y_3)$ on the $xy$ plane, where $a_0,a_1,a_2$ are constants and $x_1,x_2,x_3$ are all distinct.

  \item The curve in (3) can be represented in the form
        \[
          y=c_1y_1+c_2y_2+c_3y_3,
        \]
        where each of $c_1,c_2,c_3$ does not depend on $y_1,y_2,y_3$. Find $c_1,c_2,c_3$.

  \item Let us represent a curve
        \[
          y=a_0+a_1x+a_2x^2+a_3x^3+a_4x^4
        \]
        passing through five points $(x_1,y_1),\ldots,(x_5,y_5)$ on the $xy$ plane in the form
        \[
          y=c_1y_1+\cdots+c_5y_5,
        \]
        where each of $c_1,\ldots,c_5$ does not depend on $y_1,\ldots,y_5$, and $x_1,\ldots,x_5$ are all distinct. Find $c_1$.
\end{enumerate}

\section*{Problem 2}

Let $t$ be a real independent variable, and let $x(t)$ and $y(t)$ be real-valued functions. Answer the following questions.

\begin{enumerate}
  \item Find all solutions of the following ordinary differential equation
        \[
          \frac{d^2x}{dt^2}+2\frac{dx}{dt}+x=\cos(t),
        \]
        which are bounded when $t\to-\infty$.

  \item Find all solutions $x(t)$ and $y(t)$ of the following ordinary differential equations
        \[
          \begin{aligned}
            \frac{d^2x}{dt^2}+2\frac{dx}{dt}+x-y&=\cos(t),\\
            \frac{d^2y}{dt^2}+2\frac{dy}{dt}+y-x&=0,
          \end{aligned}
        \]
        which are bounded when $t\to-\infty$.

  \item By converting the following ordinary differential equation
        \[
          e^{-t}x^2-2\frac{dx}{dt}+x=0
        \]
        to a linear ordinary differential equation with an appropriate change of variable, find the solution $x(t)$ that satisfies $x(0)=\dfrac12$.
\end{enumerate}

\section*{Problem 3}

Let us randomly place circle stones and square stones one by one in a line from left to right. The circle and square stones are placed with probabilities $1-q$ and $q$, respectively, according to the independent and identical distribution, where $0<q<1$. The placement stops right after $M$ square stones are placed in a row, where $M$ is a positive integer. We show examples of the lines for $M=4$ as follows:
\[
  \text{line 1: }\circ\square\square\square\square,
  \qquad
  \text{line 2: }\square\circ\square\circ\circ\square\square\square\square.
\]
Let $L$ be a random variable which represents the number of the stones after stopping the placement. For the case of the lines shown above, $L=5$ and $L=9$ for lines 1 and 2, respectively.

Here, we consider intermediate states during the placement. Let $k$ be a non-negative integer and let $C_k$ be a state of a line where there are $k$ square stones in a row from the right end. For instance, we consider the following lines for $M=4$:
\[
  \text{line 3: }\circ\square\square\circ\circ\square\square,
  \qquad
  \text{line 4: }\circ\square\circ\circ.
\]
Since we are considering the case of $M=4$, lines 3 and 4 are not stopped yet. Line 3 is in state $C_2$ since there are 2 square stones in a row from the right end. Line 4 is in state $C_0$ since there is no square stone at the right end. Let $a_{kn}$ be the probability that the stopping condition is met after placing $n$ stones starting from state $C_k$, where $n$ is a non-negative integer. We define the following generating function $A_k(t)$ for $a_{kn}$:
\[
  A_k(t)=\sum_{n=0}^{\infty}t^na_{kn}.
\]

Answer the following questions.

\begin{enumerate}
  \item Calculate the mean and variance of $L$ for $M=1$.

  \item Obtain the recurrence relation that $A_k(t)$ satisfies.

  \item Obtain $A_k(t)$ as a function of $q,M,t,$ and $k$.

  \item Calculate the mean of $L$.
\end{enumerate}

\newpage

\begin{center}
  {\Large \textbf{Solutions 1--3}}
\end{center}

\section*{Solution 1}

\begin{enumerate}
  \item Expanding the determinant along the third row gives
        \[
          f(x,y)
          =
          (x_1y_2-x_2y_1)
          -(y_2-y_1)x
          +(x_2-x_1)y.
        \]
        Thus $f(x,y)=0$ is a linear equation in $x$ and $y$. Since $x_2-x_1\neq 0$, this equation is a non-degenerate line.

        Moreover,
        \[
          f(x_1,y_1)=0,
          \qquad
          f(x_2,y_2)=0,
        \]
        because in each determinant two rows are equal. Therefore the solution set is the line passing through $(x_1,y_1)$ and $(x_2,y_2)$.

  \item This is a Vandermonde determinant:
        \[
          \begin{vmatrix}
            1 & x_1 & x_1^2\\
            1 & x_2 & x_2^2\\
            1 & x_3 & x_3^2
          \end{vmatrix}
          =
          \boxed{(x_2-x_1)(x_3-x_1)(x_3-x_2).}
        \]

  \item The curve must satisfy
        \[
          \begin{bmatrix}
            1 & x_1 & x_1^2\\
            1 & x_2 & x_2^2\\
            1 & x_3 & x_3^2
          \end{bmatrix}
          \begin{bmatrix}
            a_0\\a_1\\a_2
          \end{bmatrix}
          =
          \begin{bmatrix}
            y_1\\y_2\\y_3
          \end{bmatrix}.
        \]
        The determinant of the coefficient matrix is
        \[
          (x_2-x_1)(x_3-x_1)(x_3-x_2),
        \]
        which is nonzero because $x_1,x_2,x_3$ are all distinct. Hence the coefficients $a_0,a_1,a_2$ exist uniquely.

  \item By Lagrange interpolation,
        \[
          y=c_1y_1+c_2y_2+c_3y_3,
        \]
        where
        \[
          \boxed{
          c_1=
          \frac{(x-x_2)(x-x_3)}
          {(x_1-x_2)(x_1-x_3)},}
        \]
        \[
          \boxed{
          c_2=
          \frac{(x-x_1)(x-x_3)}
          {(x_2-x_1)(x_2-x_3)},}
        \]
        and
        \[
          \boxed{
          c_3=
          \frac{(x-x_1)(x-x_2)}
          {(x_3-x_1)(x_3-x_2)}.}
        \]

  \item Again by Lagrange interpolation, the coefficient of $y_1$ in the unique quartic interpolation polynomial is
        \[
          \boxed{
          c_1=
          \frac{(x-x_2)(x-x_3)(x-x_4)(x-x_5)}
          {(x_1-x_2)(x_1-x_3)(x_1-x_4)(x_1-x_5)}.}
        \]
\end{enumerate}

\section*{Solution 2}

\begin{enumerate}
  \item The homogeneous equation is
        \[
          x''+2x'+x=0,
        \]
        whose characteristic equation is
        \[
          (r+1)^2=0.
        \]
        Hence the homogeneous solution is
        \[
          (C_1+C_2t)e^{-t}.
        \]
        This is unbounded as $t\to-\infty$ unless $C_1=C_2=0$.

        For a particular solution, put
        \[
          x_p=A\cos t+B\sin t.
        \]
        Substitution gives
        \[
          x_p''+2x_p'+x_p
          =
          2B\cos t-2A\sin t.
        \]
        Therefore $A=0$ and $B=\dfrac12$. The bounded solution is
        \[
          \boxed{x(t)=\frac12\sin t.}
        \]

  \item Put
        \[
          u=x+y,\qquad v=x-y.
        \]
        Adding and subtracting the two equations gives
        \[
          u''+2u'=\cos t,
        \]
        and
        \[
          v''+2v'+2v=\cos t.
        \]

        For the first equation, set $w=u'$. Then
        \[
          w'+2w=\cos t.
        \]
        The bounded solution as $t\to-\infty$ is
        \[
          w=\frac25\cos t+\frac15\sin t.
        \]
        Hence
        \[
          u=C+\frac25\sin t-\frac15\cos t,
        \]
        where $C$ is an arbitrary constant.

        For the second equation, the homogeneous part is unbounded as $t\to-\infty$, so only the particular solution remains:
        \[
          v=\frac15\cos t+\frac25\sin t.
        \]
        Since
        \[
          x=\frac{u+v}{2},
          \qquad
          y=\frac{u-v}{2},
        \]
        all bounded solutions are
        \[
          \boxed{
          x(t)=c+\frac25\sin t,\qquad
          y(t)=c-\frac15\cos t,}
        \]
        where $c\in\mathbb{R}$ is arbitrary.

  \item The equation can be written as
        \[
          2x'=e^{-t}x^2+x.
        \]
        Since $x(0)=1/2$, we may put
        \[
          z=\frac1x.
        \]
        Then
        \[
          z'=-\frac{x'}{x^2}.
        \]
        Therefore
        \[
          z'
          =
          -\frac12z-\frac12e^{-t},
        \]
        or
        \[
          z'+\frac12z=-\frac12e^{-t}.
        \]
        Multiplying by $e^{t/2}$ gives
        \[
          \frac{d}{dt}\left(e^{t/2}z\right)
          =
          -\frac12e^{-t/2}.
        \]
        Hence
        \[
          e^{t/2}z=e^{-t/2}+C,
        \]
        so
        \[
          z=e^{-t}+Ce^{-t/2}.
        \]
        Using $x(0)=1/2$, we get $z(0)=2$, so $C=1$. Thus
        \[
          \boxed{
          x(t)=\frac{1}{e^{-t}+e^{-t/2}}
          =
          \frac{e^t}{1+e^{t/2}}.}
        \]
\end{enumerate}

\section*{Solution 3}

\begin{enumerate}
  \item When $M=1$, the process stops at the first square stone. Thus
        \[
          P(L=n)=(1-q)^{n-1}q
          \qquad (n=1,2,\ldots).
        \]
        Hence $L$ has a geometric distribution with success probability $q$. Therefore
        \[
          \boxed{E[L]=\frac1q,}
          \qquad
          \boxed{\operatorname{Var}(L)=\frac{1-q}{q^2}.}
        \]

  \item For $0\leq k\leq M-2$, after placing one stone, a circle stone brings the state to $C_0$, while a square stone brings the state to $C_{k+1}$. Hence
        \[
          \boxed{
          A_k(t)=t\{(1-q)A_0(t)+qA_{k+1}(t)\}
          \qquad (0\leq k\leq M-2).}
        \]
        For $k=M-1$, a square stone stops the placement. Therefore
        \[
          \boxed{
          A_{M-1}(t)=t\{(1-q)A_0(t)+q\}.}
        \]

  \item Put
        \[
          z=tq.
        \]
        From the recurrence in (2), backward substitution gives
        \[
          A_k(t)
          =
          t(1-q)A_0(t)(1+z+\cdots+z^{M-k-1})
          +z^{M-k}.
        \]
        Setting $k=0$ and solving for $A_0(t)$ gives
        \[
          A_0(t)
          =
          \frac{(tq)^M(1-tq)}
          {1-t+t(1-q)(tq)^M}.
        \]
        Substituting this into the expression for $A_k(t)$ yields
        \[
          \boxed{
          A_k(t)=
          \frac{(tq)^{M-k}(1-t)+t(1-q)(tq)^M}
          {1-t+t(1-q)(tq)^M}}
        \]
        for $k=0,1,\ldots,M-1$.

  \item Since the process starts from state $C_0$, the generating function of $L$ is $A_0(t)$. Therefore
        \[
          E[L]=A_0'(1).
        \]
        Using
        \[
          A_0(t)=
          \frac{(tq)^M(1-tq)}
          {1-t+t(1-q)(tq)^M},
        \]
        logarithmic differentiation at $t=1$ gives
        \[
          A_0'(1)
          =
          \frac{1-q^M}{(1-q)q^M}.
        \]
        Hence
        \[
          \boxed{
          E[L]=\frac{1-q^M}{(1-q)q^M}.}
        \]
\end{enumerate}

\end{document}
